
%% OLD VERSION
% Machine Unlearning (MU) presents a critical challenge for Large Language Models (LLMs), yet current methods overlook a key factor: \textit{a model's pre-existing familiarity with the knowledge to be forgotten}. We hypothesize that popular facts are more deeply embedded in a model's parameters than rare ones, posing a distinct unlearning challenge. Previously, \textit{most unlearning benchmarks used synthetic data, failing to account for popularity and lacking realism} by not including data potentially used in pre-training. To address this, we introduce \textbf{UnlearnPopQA} -- the first benchmark, based on the PopQA dataset, designed to evaluate unlearning across the spectrum of fact popularity on models of various sizes. It uses real-world Wikidata, which could have been part of the pre-training corpus.
% Our experiments reveal a \textit{interplay between knowledge popularity and model scale}. For smaller models (1B, 3B), we observe little difference in unlearning rare versus popular facts. However, as model scale increases, a clear dependency emerges: for large LLMs, attempting to unlearn popular knowledge is significantly more likely to induce catastrophic forgetting than unlearning rare knowledge. Despite this volatility in factual knowledge on UnlearnPopQA, the models' general reasoning capabilities, measured on benchmarks like MMLU and HellaSwag, remain surprisingly robust.
% Our work provides that for large-scale LLMs, a \textit{one-size-fits-all unlearning strategy is insufficient}. We conclude that it's critical to develop popularity-aware approaches that carefully calibrate unlearning parameters for rare and popular entities to achieve effective and safe knowledge removal.



%% NEW VERSION
Machine Unlearning (MU) is an increasingly critical capability for Large Language Models (LLMs), enabling the removal of unsafe, outdated, or private information. However, existing unlearning methods are often evaluated under the assumption that all facts are equally challenging to forget. In this paper, we approach MU from a new angle: \textit{Does the popularity of a fact influence how effectively an LLMs can unlearn it?}

To answer this question, we propose \textbf{UNLamb}, the benchmark designed to systematically investigate this relationship. UNLamb consists of 11.6k question-answer pairs derived from real-world knowledge in Wikidata, explicitly partitioned into rare and popular facts based on entity prominence. Using this benchmark, we perform a comprehensive evaluation of state-of-the-art unlearning algorithms on a set of LLMs of different sizes.
We conduct a comprehensive analysis of four unlearning methods: GA, GD, NPO, and RMU, across three validation sets and two LLMs benchmarks. Our results show that larger models struggle more to forget popular facts and often damage related knowledge when trying to remove them. In contrast, it is much easier to remove rare facts without side effects. This shows that current unlearning methods do not suffice and highlights the need for new strategies that account for the popularity of facts. 
% Our experiments reveal a crucial interplay between model scale and fact popularity. For large-scale models, we show that unlearning popular facts is significantly more difficult and more likely to cause collateral damage to related knowledge than unlearning rare facts. Our work demonstrates that a classic approach to unlearning is insufficient, highlighting the critical need for popularity-aware strategies to ensure safe and effective knowledge removal in modern LLMs.